\section{总结与延伸}

本节把 EP110 压缩成几个结论。第一，Agent 的最小定义是感知和行动，而不是聊天。第二，Kimi K2、ChatGPT Agent、Qwen3-Coder 和 Manus 分别展示了 Agent 栈的模型、产品、代码/RL 和上下文工程层。第三，Agent training 的核心从 input-output 走向 trajectory、environment、task-reward 和 safety。第四，生产 Agent 的关键是系统工程：KV cache、工具遮蔽、文件系统记忆、错误恢复和用户控制。

\begin{importantbox}{把 EP110 放进张小珺 AI 队列}
EP115 给出 Agent 下半场理论，EP113 讲 Kimi K2 的模型公司视角，EP116 讲企业级 Agentic Model；EP110 则把几份最新技术报告合在一起，展示 Agent 从论文、模型、产品到工程落地的完整链条。
\end{importantbox}

\begin{knowledgebox}{与前后几集的关系}
\begin{tabularx}{\textwidth}{p{0.18\textwidth}p{0.34\textwidth}X}
\toprule
节目 & 主题 & 与 EP110 的连接 \\
\midrule
EP115 & Agent 下半场理论 & 给出 reward、environment、interface 的理论框架。 \\
EP113 & Kimi K2 和 Agentic LLM & 给出模型公司如何理解 K2 和开源生态。 \\
EP116 & 企业级 Agentic Model & 给出 Agent 在 ToB 私有数据中的部署视角。 \\
EP139 & Agent 技术史 & 给出 Agent 从早期系统到 LLM Agent 的历史脉络。 \\
\bottomrule
\end{tabularx}
\end{knowledgebox}

\subsection{关键 takeaways}

前面章节已经分别从定义、训练、产品、代码环境和上下文工程解释 Agent。本节把这些内容压缩成几条工程判断，方便后续和 EP115、EP113、EP116 对照。

\begin{enumerate}
\item Agent = 感知 + 行动 + 反馈，不是“会调用工具”的单点能力。
\item Kimi K2 的重点是开放 MoE 模型、MuonClip、合成工具轨迹和 joint RL。
\item ChatGPT Agent 的重点是统一产品系统和用户控制。
\item Qwen3-Coder 的重点是代码环境、Code RL、Long-horizon RL 和工具链。
\item Manus 的重点是上下文工程，特别是 KV cache、工具遮蔽和文件系统记忆。
\end{enumerate}

\subsection{开放问题}

这些问题是 Agent 从报告走向产品后仍然没有定论的部分。它们决定下一轮训练范式、产品架构和基础设施投入方向。

\begin{enumerate}
\item Agent 训练中，合成轨迹和真实环境反馈如何配比？
\item Rubrics as reward 能否支撑开放任务中的可靠 self-improvement？
\item In-context 与 end-to-end 两条路线最终会融合到什么形态？
\item 生产 Agent 的成本瓶颈会更主要来自模型调用、工具执行还是上下文长度？
\end{enumerate}

\subsection{拓展阅读}

\begin{itemize}
\item 对 Agent 理论框架感兴趣，可对照 EP115 姚顺雨访谈。
\item 对 Kimi K2 模型公司视角感兴趣，可对照 EP113 杨植麟访谈。
\item 对企业 Agentic Model 感兴趣，可对照 EP116 吴明辉访谈。
\end{itemize}

\end{document}
